\begin{enumerate}[label=\textbf{\ansnum.\arabic*}]
  \item Copy: \verb!i64!, \verb!(f64, f64)!, \verb![i32; 4]!, \verb!&str! (a shared reference is a small value that may be copied freely), \verb!bool!. Moved: \verb!String!, \verb!Vec<u8>! (they own heap memory).
  \item \verb!let t = s;! moves the ownership of the string to \verb!t!; \verb!s! is then invalid, and \verb|println!("{s}")| uses it (rule 2: one owner at a time). Fixes: \verb!let t = s.clone();! (two independent strings), or \verb!let t = &s;! (borrow; \verb!t! is a reference and \verb!s! remains the owner). Printing \verb!t! instead of \verb!s! also works but loses \verb!s!.
  \item
%rs:a3
\begin{verbatim}
fn count_above(xs: &[f64], limit: f64) -> usize {
    xs.iter().filter(|&&x| x > limit).count()
}

fn main() {
    let v = vec![0.2, 1.5, 0.9, 2.2, 3.1, 0.1];
    println!("{}", count_above(&v, 1.0));
    println!("{}", count_above(&v[3..], 1.0));
    println!("{:?}", v);
}
\end{verbatim}
prints $3$, $2$ and the unchanged vector. The function only \emph{borrows}: it receives a read-only view and the ownership stays in \verb!main!.
  \item
%rs:a4
\begin{verbatim}
fn clamp_all(xs: &mut [f64], lo: f64, hi: f64) {
    for x in xs.iter_mut() {
        *x = x.clamp(lo, hi);
    }
}

fn main() {
    let mut v = vec![-2.0, 0.5, 3.0];
    clamp_all(&mut v, -1.0, 1.0);
    println!("{:?}", v);
}
\end{verbatim}
prints \verb![-1.0, 0.5, 1.0]!.
  \item The line \verb!v.push(4);! needs a mutable borrow while \verb!r! (a shared borrow) is still used afterwards in \verb|println!|. Fix 1: use \verb!r! \emph{before} the push (print first, push after: the shared borrow has ended). Fix 2: do not borrow, but copy: \verb!let r = v.clone();! so \verb!r! is independent of \verb!v!.
  \item \verb!process(f: Frame)! takes the ownership: the caller cannot use the frame afterwards. \verb!&Frame! only reads, \verb!&mut Frame! may change the frame but returns it afterwards. Mean colour: \verb!&Frame!. Blur in place: \verb!&mut Frame!. Send to the logging thread which stores it: \verb!Frame! (ownership moves into the thread, which then decides when it is dropped; no copy and no sharing problems).
  \item A \verb!&str! parameter accepts a \verb!&String! (it is converted automatically to a view of its text) and a literal, which is already a \verb!&str!. A \verb!String! parameter must \emph{own} its text, so a literal does not fit. The caller converts: \verb!"abc".to_string()!, \verb!String::from("abc")! or \verb!"abc".to_owned()! (each allocates and copies).
  \item A reference that points to memory that has already been freed. Example:
%rsfail:a8
\begin{verbatim}
fn make() -> &String {
    let s = String::from("goal");
    &s
}

fn main() {
    println!("{}", make());
}
\end{verbatim}
The compiler reports \verb!error[E0106]: missing lifetime specifier! with the hint that the function returns a borrowed value but there is nothing it could borrow from; even with a lifetime written out it would say that \verb!s! does not live long enough, because it is dropped at the end of \verb!make!. Fix: return the owned value, \verb!fn make() -> String { String::from("goal") }!.
  \item A clone copies $10^7\times8=80$\,MB each time: slow (milliseconds to tens of milliseconds) and memory-hungry. Passing by value as a move copies only the handle (24 bytes), almost free, but the caller loses the vector and could not use it in the next iteration. A reference is also almost free ($16$ bytes) and the caller keeps the data. For a function that only reads, pass \verb!&[f64]!.
  \item
%rs:a10
\begin{verbatim}
fn split_pair(xs: &mut [f64]) -> (&mut [f64], &mut [f64]) {
    let mid = xs.len() / 2;
    xs.split_at_mut(mid)
}

fn main() {
    let mut v = vec![1.0, 2.0, 3.0, 4.0, 5.0, 6.0];
    let (a, b) = split_pair(&mut v);
    for x in a.iter_mut() { *x = 0.0; }
    for x in b.iter_mut() { *x *= 2.0; }
    println!("{:?}", v);
}
\end{verbatim}
prints \verb![0.0, 0.0, 0.0, 8.0, 10.0, 12.0]!. The rule protects against \emph{two names for the same memory} changing it: the two halves do not overlap, so they cannot interfere. The function \verb!split_at_mut! is written (with a small amount of checked, hand-verified unsafe code inside the standard library) to promise exactly that; for you it is an ordinary safe function.
\end{enumerate}
